\section{Selected Realizations of Relational Fidelity}
\label{sec:relational-distortion-realizations}

Relational requirements may be supplied directly by a source relation or emerge from reconstruction.
This section instantiates both routes within the finite-codeword geometry, covering signed and baseline-relative pair requirements, graph-operator probes, and exact correspondences to centroid reconstruction and normalized cut.
These examples are representative rather than exhaustive; path, walk, and additional spectral constructions are retained in \ref{app:path-walk-content}.

\subsection{Selected constructions for relational requirements}

When pairwise weights are derived from a single real-valued relation, two useful cases are an intrinsically signed relation and a relation interpreted relative to a source-side baseline.

This role assignment is specific to the collision realization: for example, the block-valued decoder of \cref{sec:block-summary-walkthrough} can reproduce strong relations between different groups through its reproduction table $\Xi$, without interpreting group inequality itself as a separation requirement.

\paragraph{Intrinsically signed relations}
For an intrinsically signed $S$, the nonnegative source weights can be $S_+=[S]_+$ and $S_-=[-S]_+$.

\paragraph{Baseline-relative relations}
Alternatively, compare a real-valued relation with a source-side neutral baseline:
\begin{equation}
S_+=[S-S_{\mathrm{bg}}]_+,\qquad
S_-=[S_{\mathrm{bg}}-S]_+.
\label{eq:signed-relation}
\end{equation}
Relations above the baseline induce alignment and those below it induce separation, including when all original source values are nonnegative.
When a probability-law form is desired, nonzero pair weights can be normalized over the evaluated carrier pairs to obtain $P_+$ and $P_-$; weighted-sum objectives may instead use them directly.

The source-side baseline $S_{\mathrm{bg}}$ identifies unusually strong or weak source relationships, while $r(z)$ specifies a reference for codeword occupancy.
A representation-side collision reference can separately center a learning margin.
These are respectively a source-side relational baseline, a codeword-occupancy reference, and a representation-side collision threshold.

\subsection{Direct relation preservation and signed distinctions}
\label{sec:edge-cut-distortion}

The simplest graph realization treats retained edge weight itself as relational content and then uses collision to decide which source edges survive.

Let $G=(V,E,W)$ be an undirected graph with symmetric nonnegative weights, zero diagonal, and one copy of each undirected edge in $E$, and let $\widetilde G$ denote a reconstructed weighted graph on the same carrier with weight matrix $\widetilde W$.
A reconstruction obtained by deletion or attenuation satisfies $0\le\widetilde W\le W$ entrywise.
The direct edge content and its distortion are
\begin{equation}
\mathcal I_{\mathrm{edge}}(G)=\frac12\sum_{i,j}W_{ij},\qquad
D_{\mathrm{edge}}(G,\widetilde G)=\frac12\sum_{i,j}(W_{ij}-\widetilde W_{ij}).
\label{eq:edge-distortion}
\end{equation}
When $\mathcal I_{\mathrm{edge}}(G)>0$, the normalized direct-edge distortion is
\begin{equation}
D_E(G,\widetilde G)
:=\frac{D_{\mathrm{edge}}(G,\widetilde G)}{\mathcal I_{\mathrm{edge}}(G)}.
\label{eq:normalized-edge-distortion}
\end{equation}
The factor $1/2$ counts each undirected edge once.
This is a monotone content-based fidelity in the sense of \cref{eq:relational-content-distortion}.

A hard partition can retain only within-codeword edges,
\begin{equation}
\widetilde W_{ij}=W_{ij}\ind[c(i)=c(j)],
\label{eq:partition-retained-relation}
\end{equation}
so its distortion is the cut weight $\sum_{\{i,j\}\in E}W_{ij}\ind[c(i)\ne c(j)]$.
For probabilistic assignments, define $q_i:=q_\theta(\cdot\mid x_i)$ and $k_{\theta,ij}:=k_\theta(x_i,x_j)=q_i^\top q_j$.
Collision attenuation then gives
\begin{equation}
\widetilde W_{\theta,ij}=W_{ij}k_{\theta,ij},\qquad
D_{\mathrm{edge}}(\theta)=\frac12\sum_{i,j}W_{ij}(1-k_{\theta,ij}).
\label{eq:soft-cut-distortion}
\end{equation}
Its normalized form is
\[
D_E(\theta)
=\frac{
\sum_{\{i,j\}\in E}W_{ij}(1-k_{\theta,ij})
}{
\sum_{\{i,j\}\in E}W_{ij}
}.
\]
When used as a learning objective, minimizing this distortion is equivalent to maximizing collision on positively weighted source edges.
For independent sampled assignments, the soft expression also equals the expected hard cut weight on distinct vertices because this distortion is linear in the retained edge indicators; this equivalence does not generally extend to nonlinear graph fidelities.

The decoder convention matters here.
A reconstruction of the retained weighted graph includes its weights as stored or shared information.
An equality-only description can instead be evaluated by its cut distortion against the source graph.
These distinct descriptions can lead to the same scalar distortion.

Positive edge preservation alone allows complete collapse.
For nonnegative source weights $W_+$ and $W_-$ describing favored and disfavored collisions, a signed relational objective is
\begin{equation}
\mathcal L_{\mathrm{signed}}(\theta)
=\frac12\sum_{i,j}\left[W_{+,ij}(1-k_{\theta,ij})
+\lambda_{\mathrm{rel}}W_{-,ij}k_{\theta,ij}\right].
\label{eq:signed-relational-objective}
\end{equation}
The weights may be supplied directly or derived from $W-W_{\mathrm{bg}}$ as in \cref{eq:signed-relation}.
Although collision attenuation still gives $0\le\widetilde W\le W$ for a nonnegative source graph, the signed objective is not a monotone retained-content distortion: decreasing collision worsens fidelity on favored pairs but improves it on disfavored pairs.
With a hard equality decoder, it defines a relational fidelity criterion on the partition relation, closely related to the positive/negative disagreement perspective of correlation clustering~\citep{bansal2004correlation}.

\subsection{Graph probes and operator fidelity}
\label{sec:graph-fourier-distortion}

Direct edge fidelity asks which relations remain; operator fidelity instead asks which behaviors supported by those relations remain.

For the same undirected graph, let $d_i=\sum_jW_{ij}$ and $L=\operatorname{Diag}(d)-W$.
Its Dirichlet energy is
\begin{equation}
f^\top Lf=\frac12\sum_{i,j}W_{ij}(f_i-f_j)^2.
\label{eq:dirichlet-energy-vertex}
\end{equation}
With $L=U\Lambda U^\top$ and graph-Fourier coefficients $\widehat f=U^\top f$, the same scalar is
\begin{equation}
f^\top Lf=\sum_r\lambda_r|\widehat f_r|^2.
\label{eq:dirichlet-energy-fourier}
\end{equation}
This is the standard graph-signal/Fourier interpretation~\citep{shuman2013emerging}.
The relation determines which signal variations carry energy.

Deleting or attenuating edges gives $0\preceq\widetilde L\preceq L$, since $L-\widetilde L$ is the Laplacian of the removed weights.
A fixed probe therefore loses
\begin{equation}
D_f(G,\widetilde G)
=f^\top(L-\widetilde L)f
=\frac12\sum_{i,j}(W_{ij}-\widetilde W_{ij})(f_i-f_j)^2.
\label{eq:dirichlet-distortion-vertex}
\end{equation}
If $\widetilde L=\widetilde U\widetilde\Lambda\widetilde U^\top$, this is equally
\begin{equation}
D_f(G,\widetilde G)
=\sum_r\lambda_r|(U^\top f)_r|^2
-\sum_r\widetilde\lambda_r|(\widetilde U^\top f)_r|^2.
\label{eq:dirichlet-distortion-fourier}
\end{equation}
The vertex and Fourier expressions evaluate the same fidelity in two coordinate systems, using each operator's own eigenbasis.

Rather than selecting a single probe signal, we can specify a distribution over signals whose behavior we wish to preserve.
Let $F$ be such a random probe with second moment $M=\E[FF^\top]$.
Specializing \cref{eq:operator-expected-distortion} to the graph Laplacian gives
\[
D_M(G,\widetilde G)=\operatorname{tr}[(L-\widetilde L)M].
\]
For an edge $e=\{i,j\}$, let $b_e=e_i-e_j$ be an oriented incidence vector, with either orientation.
The Laplacian decomposes into individual edge contributions:
\[
L
=
\sum_{e=\{i,j\}\in E}W_{ij}b_eb_e^\top.
\]
The orientation is arbitrary because replacing $b_e$ by $-b_e$ leaves $b_eb_e^\top$ unchanged.
Using the rank-one trace identity $\operatorname{tr}(uv^\top A)=v^\top Au$ gives
\begin{align*}
\operatorname{tr}(LM)
&=
\sum_{e=\{i,j\}\in E}
W_{ij}\operatorname{tr}(b_eb_e^\top M)\\
&=
\sum_{e=\{i,j\}\in E}
W_{ij}b_e^\top M b_e.
\end{align*}
Because $M=\E[FF^\top]$,
\[
b_e^\top M b_e
=
\E[(F_i-F_j)^2].
\]
Define the source importance of edge $e$ by $s_e=W_{ij}b_e^\top M b_e\ge0$.
It is the contribution of that edge to the expected source Dirichlet energy under the chosen probe family.
Since $\sum_e s_e=\operatorname{tr}(LM)$, collision attenuation gives the normalized distortion
\begin{equation}
\overline D_M(\theta)
=\sum_{e=\{i,j\}\in E}\rho_e(1-k_{\theta,ij}),\qquad
\rho_e=\frac{s_e}{\sum_{e'\in E}s_{e'}},
\label{eq:probe-edge-family}
\end{equation}
provided the total expected source energy $\operatorname{tr}(LM)$ is positive.
Thus $M$ determines which source signal variations, and therefore which edges, matter most to the fidelity criterion, while the finite-code representation itself remains unchanged.

\paragraph{Isotropic and all-mode probes}
For $M=I$, $D_I=\operatorname{tr}(L-\widetilde L)=2D_{\mathrm{edge}}$.
Since $\operatorname{tr}(L)=2\mathcal I_{\mathrm{edge}}(G)$, the normalized isotropic-probe criterion is $D_I/\operatorname{tr}(L)=D_E$.
For a source graph with $n$ vertices and $\kappa_G$ connected components, suppose $\operatorname{rank}(L)=n-\kappa_G>0$.
A source-normalized alternative then uses $M=L^+/(n-\kappa_G)$, where $L^+$ is the Moore--Penrose pseudoinverse.
We denote this source-normalized all-mode criterion by $D_F$.
The spectral expansion
\[
L^+
=
\sum_{\lambda_r>0}\frac1{\lambda_r}u_ru_r^\top
\]
gives
\begin{align}
D_F(G,\widetilde G)
&=\frac{\operatorname{tr}[(L-\widetilde L)L^+]}{n-\kappa_G}\\
&=\frac1{n-\kappa_G}
\sum_{\lambda_r>0}
\frac{u_r^\top(L-\widetilde L)u_r}{\lambda_r}.
\label{eq:source-normalized-dirichlet-distortion}
\end{align}
The factor $1/\lambda_r$ normalizes each positive-eigenvalue source Laplacian mode by its source Dirichlet energy, while division by $n-\kappa_G$ averages over those modes.
Zero modes are excluded because they have zero source Dirichlet energy; for a disconnected graph, they include signals that are constant separately on each connected component.
For an individual positive-eigenvalue source mode $u_r$, its source Dirichlet energy is $u_r^\top Lu_r=\lambda_r$.
Evaluating the same fixed source signal on the reconstructed graph gives the Dirichlet energy $u_r^\top\widetilde Lu_r$, so
\[
\frac{u_r^\top(L-\widetilde L)u_r}{\lambda_r}
=1-\frac{u_r^\top\widetilde Lu_r}{\lambda_r}
\]
is the fraction of that mode's source Dirichlet energy removed by compression, and $D_F$ is the mean of these relative losses.
The quantity $u_r^\top\widetilde Lu_r$ is not generally an eigenvalue of $\widetilde L$, because a source eigenvector need not remain an eigenvector after compression.
Within a repeated positive-eigenvalue source eigenspace, the individual mode losses depend on the chosen orthonormal eigenbasis, but their sum over the entire eigenspace and hence its contribution to $D_F$ is basis invariant.
Because $\operatorname{tr}(LL^+)=\operatorname{rank}(L)=n-\kappa_G$, the normalization $M=L^+/(n-\kappa_G)$ ensures $\operatorname{tr}(LM)=1$, so the probe distribution has unit expected Dirichlet energy on the source graph.

The spectral form describes distortion mode by mode.
To obtain the complementary edgewise interpretation, decompose the removed Laplacian into its deleted or attenuated edge contributions:
\[
L-\widetilde L
=
\sum_{e=\{i,j\}\in E}(W_{ij}-\widetilde W_{ij})b_e b_e^\top.
\]
Substituting this decomposition into the trace form gives
\begin{equation}
D_F(G,\widetilde G)
=\frac1{n-\kappa_G}\sum_{e=\{i,j\}\in E}(W_{ij}-\widetilde W_{ij})R_{\mathrm{eff}}^G(i,j),
\qquad R_{\mathrm{eff}}^G(i,j)=b_e^\top L^+b_e,\quad e=\{i,j\}.
\label{eq:resistance-weighted-dirichlet-distortion}
\end{equation}
All resistances are computed in the original source graph.
For a hard partition,
\begin{equation}
D_F(G,c)=\frac1{n-\kappa_G}
\sum_{\substack{\{i,j\}\in E\\c(i)\ne c(j)}}
W_{ij}R_{\mathrm{eff}}^G(i,j).
\label{eq:hard-partition-dirichlet-distortion}
\end{equation}
Its numerator is the established resistance-weighted cut-size~\citep{cesabianchi2013random,gentile2013online}.
Here it is used as the all-mode source-normalized probe realization of relational fidelity.
Since
\[
\sum_{\{i,j\}\in E}W_{ij}R_{\mathrm{eff}}^G(i,j)
=
\operatorname{tr}(LL^+)
=
n-\kappa_G,
\]
the denominator is the total source resistance-weighted edge mass, giving $0\le D_F\le1$ under attenuation.

\paragraph{Source collision probes}
Beyond isotropic and all-mode choices, a source-derived probe family can emphasize distinctions already present in the graph's local transition structure.
Viewing each vertex's normalized edge weights as a distribution over destinations, two vertices have greater probability-product overlap when independent draws from their destination distributions tend to coincide.
For positive degrees, let $P_{ij}=W_{ij}/d_i$ be the source transition matrix.
Its rows are distributions over relational destinations, and
\begin{equation}
\Gamma_{\mathrm{src}}=PP^\top,\qquad
(\Gamma_{\mathrm{src}})_{ij}=\sum_rP_{ir}P_{jr}
\end{equation}
is the probability that independent destination draws from rows $i$ and $j$ coincide.
It is positive semidefinite and can be used as a source-frozen probe second moment.
For $\operatorname{tr}(L\Gamma_{\mathrm{src}})>0$,
\begin{equation}
D_C(G,\widetilde G)
=\frac{\operatorname{tr}[(L-\widetilde L)\Gamma_{\mathrm{src}}]}
{\operatorname{tr}(L\Gamma_{\mathrm{src}})}.
\label{eq:source-collision-covariance-distortion}
\end{equation}
For $e=\{i,j\}$, $b_e^\top PP^\top b_e=\|P_i-P_j\|_2^2$, so collision attenuation gives
\begin{equation}
D_C(\theta)=
\frac{\sum_{\{i,j\}\in E}W_{ij}\|P_i-P_j\|_2^2(1-k_{\theta,ij})}
{\sum_{\{i,j\}\in E}W_{ij}\|P_i-P_j\|_2^2}.
\label{eq:source-collision-edge-distortion}
\end{equation}
This fidelity emphasizes differences in source transition profiles across edges.
Here $\Gamma_{\mathrm{src}}=PP^\top$ is a source-derived probe second moment, encoding collision probabilities between transition distributions, whereas $k_{\theta,ij}=q_i^\top q_j$ is a representation-side collision probability that controls edge attenuation.
The two have the same probability-product form but act on different distributions and play different roles.

The full transition-profile construction compares the destination distributions themselves.
A complementary source-derived probe can instead summarize each vertex's local transition structure by a scalar measure of its effective relational diversity, and preserve how that quantity varies across the graph.
Specifically, a rank-one source probe uses local transition collision entropy,
\begin{equation}
h_i^{(2)}=-\log\sum_jP_{ij}^2.
\end{equation}
The source field $h_G^{(2)}$ is held fixed while the graph is compressed.
When its source Dirichlet energy is positive, define
\begin{equation}
D_{H_2}(G,\widetilde G)
=\frac{(h_G^{(2)})^\top(L-\widetilde L)h_G^{(2)}}
{(h_G^{(2)})^\top Lh_G^{(2)}}.
\label{eq:entropy-probe-distortion}
\end{equation}
Here $D_{H_2}$ measures distortion of the source transition-entropy field $h_G^{(2)}$ under graph compression.
Under collision attenuation this has edge weights proportional to $W_{ij}(h_i^{(2)}-h_j^{(2)})^2$.
It measures the lost gradient energy of the selected source-intrinsic field.
If the field is constant on connected components, the denominator is zero and this normalized fidelity is undefined.

For an unweighted simple graph, each transition row is uniform over $d_i$ neighbors, so $h_i^{(2)}=\log d_i$; its Shannon row entropy is identical.
Thus on such graphs the transition-entropy field reduces to the log-degree field, while weighted or otherwise nonuniform relations allow Shannon and R\'enyi-2 transition entropies to distinguish different local structure.

\paragraph{Worst-case fidelity}
The relative worst-case graph distortion is
\begin{equation}
D_{\mathrm{wc}}(G,\widetilde G)
=\sup_{f:f^\top Lf>0}
\frac{f^\top(L-\widetilde L)f}{f^\top Lf}.
\label{eq:graph-worst-case-spectral-distortion}
\end{equation}
If a connected source becomes disconnected, a signal constant on each new component but nonconstant on the original graph has zero reconstructed energy and positive source energy.
Hence $D_{\mathrm{wc}}=1$.
In particular, every nontrivial hard partition of a connected source that removes all cross-class edges has maximal distortion under this criterion.
That makes worst-case fidelity a useful boundary of the representation, but an uninformative ranking among such hard partition reconstructions.

\subsection{Centroid reconstruction as an exact objective correspondence}
\label{sec:vq}

Classical squared-error centroid reconstruction admits an exact pairwise relational formulation.
Consider $N$ observations $x_1,\ldots,x_N\in\R^d$, and define the source relation
\[
S(i,j)=\|x_i-x_j\|_2^2.
\]
For the probabilistic finite encoder of \cref{eq:discrete-encoder}, write
\[
q_{iz}:=q_\theta(z\mid x_i),
\qquad
q_{iz}\ge0,
\qquad
\sum_zq_{iz}=1,
\]
for the soft assignment weight of observation $x_i$ to codeword $z$.
Define the empirical aggregate codeword distribution by
\[
\hat q_Z(z)
:=
\frac1N\sum_{i=1}^Nq_{iz}.
\]
This is the finite-sample analogue of the population aggregate codeword distribution $q_Z$ in \cref{eq:aggregate-code}.
A codeword is active when $\hat q_Z(z)>0$.
For each active codeword, the assignment-weighted centroid is
\begin{equation}
\bar x_z
=
\frac{1}{N\hat q_Z(z)}
\sum_{i=1}^Nq_{iz}x_i,
\qquad
\hat q_Z(z)>0.
\label{eq:finite-soft-centroid}
\end{equation}
Define the average centroid reconstruction error by
\[
D_{\mathrm{cent}}
:=
\frac1N
\sum_{z:\hat q_Z(z)>0}
\sum_iq_{iz}\|x_i-\bar x_z\|_2^2.
\]
The decoder associates one reproduction vector $\widehat x_z\in\R^d$ with each codeword $z\in\mathcal Z$.
For arbitrary choices of these reproduction vectors, define the average reconstruction error by
\[
D_{\mathrm{dec}}
:=
\frac1N
\sum_{i,z}q_{iz}\|x_i-\widehat x_z\|_2^2.
\]
The key correspondence is that centroid reconstruction error can be written exactly as an inverse-marginally weighted pairwise squared-distance distortion under the same assignments.
More generally, reconstruction with arbitrary codeword reproductions decomposes into this centroid error plus the excess incurred when the reproduction vectors differ from their assignment-weighted centroids.
The following proposition collects the established weighted variance identities that yield these correspondences~\citep{gray1998quantization,dhillon2004kernel}.

\begin{proposition}[Centroid and pairwise forms]
\label{prop:centroid-pairwise}
For the assignments and reconstruction errors above,
\begin{equation}
D_{\mathrm{cent}}
=\frac1{2N^2}
\sum_{z:\hat q_Z(z)>0}\frac1{\hat q_Z(z)}
\sum_{i,j}q_{iz}q_{jz}\|x_i-x_j\|_2^2.
\label{eq:finite-soft-pairwise}
\end{equation}
The affinity $a_{\hat q_Z}$ is the empirical finite-sample specialization of the inverse-mass affinity in \cref{eq:qz-affinity}, so equivalently
\begin{equation}
D_{\mathrm{cent}}=\frac1{2N^2}\sum_{i,j}
\|x_i-x_j\|_2^2a_{\hat q_Z}(x_i,x_j).
\label{eq:finite-soft-qz-normalized}
\end{equation}
The arbitrary-reproduction error decomposes as
\begin{equation}
D_{\mathrm{dec}}
=D_{\mathrm{cent}}
+\sum_{z:\hat q_Z(z)>0}
\hat q_Z(z)\|\bar x_z-\widehat x_z\|_2^2.
\label{eq:decoder-centroid-decomposition}
\end{equation}
\end{proposition}

\begin{proof}
For each active codeword,
\begin{equation}
\sum_iq_{iz}\|x_i-\bar x_z\|^2
=\sum_iq_{iz}\|x_i\|^2
-N\hat q_Z(z)\|\bar x_z\|^2.
\end{equation}
Expanding the weighted double sum of $\|x_i-x_j\|^2$ yields twice $N\hat q_Z(z)$ times the same within-codeword variance expression.
Summing over codewords proves the first identity, and division by the empirical aggregate mass gives the inverse-marginal affinity form.
For the decoder identity, expand $x_i-\widehat x_z=(x_i-\bar x_z)+(\bar x_z-\widehat x_z)$; the weighted cross term vanishes, and $\sum_iq_{iz}=N\hat q_Z(z)$ gives the stated excess-error term.
\end{proof}

For a hard assignment $c$, let $V_z:=\{i:c(i)=z\}$.
For each nonempty class, the familiar special case is
\begin{equation}
\sum_{i\in V_z}\|x_i-\bar x_z\|^2
=\frac1{2|V_z|}\sum_{i,j\in V_z}\|x_i-x_j\|^2,
\label{eq:vq-pairwise}
\end{equation}
The variance identity therefore determines the inverse empirical-marginal factor.
Where the assignments are differentiable in the encoder parameters and active aggregate masses remain positive, the two forms also have identical encoder gradients.

The population form requires a finite second moment.
Let $X\sim\mu$ with $\E\|X\|^2<\infty$, let $Z\mid X\sim q_\theta(\cdot\mid X)$, and define $\bar x_z=\E[X\mid Z=z]$ for active words.
Then
\begin{equation}
\E\|X-\bar x_Z\|^2
=\frac12\E_{X,X'\overset{\mathrm{iid}}{\sim}\mu}
\left[\|X-X'\|^2a_{q_Z}(X,X')\right].
\label{eq:stochastic-vq-pairwise}
\end{equation}
To see this, condition on each active $z$, apply the two-independent-copies variance identity, and average with weight $q_Z(z)$.
The conditional law has density $q_\theta(z\mid x)/q_Z(z)$ with respect to $\mu$, producing exactly the inverse-mass factor.

This correspondence concerns a scalar fidelity and its optimal per-codeword reproduction.
Pointwise reconstruction uses the reproduction vectors, while within-code relational distortion can use the partition or affinity as its reconstructed statistic.
The identity relies on this per-codeword decoder structure; a decoder that jointly interprets multiple latent variables need not admit the same decomposition.

\subsection{Normalized association and normalized cut}
\label{sec:normalized-cut-realization}

Normalized cut is a classical graph-partitioning objective that penalizes edges crossing between groups relative to the total edge volume of those groups~\citep{shi2000normalized}.
Equivalently, normalized association rewards edge weight retained within each group after normalization by group volume.
We show here that the endogenous inverse-mass affinity of the finite-codeword construction recovers these objectives under graph-local degree weighting.

For an undirected nonnegative graph with positive total volume, let $d_i=\sum_jW_{ij}$, $\operatorname{vol}(G)=\sum_i d_i$, and $\pi_i=d_i/\operatorname{vol}(G)$.
This defines a degree-weighted distribution over the vertices.
For assignments $q_i(z)$, the aggregate codeword distribution under this degree weighting is
\begin{equation}
\bar q_G(z)=\sum_i\pi_iq_i(z)
=\frac{\sum_i d_iq_i(z)}{\operatorname{vol}(G)}.
\label{eq:graph-aggregate-code}
\end{equation}
Zero-degree vertices receive zero weight under this distribution.
Write $\mathcal Z_G^+=\{z:\bar q_G(z)>0\}$.
The endogenous inverse-mass affinity $a_{q_Z}$ of \cref{eq:qz-affinity}, built from the aggregate codeword distribution $q_Z(z)$ in \cref{eq:aggregate-code}, specializes here by replacing $q_Z(z)$ with the graph-local aggregate $\bar q_G(z)$ induced by the degree-weighted vertex distribution $\pi$.
\begin{equation}
a_{\bar q_G}(i,j)
:=\sum_{z\in\mathcal Z_G^+}\frac{q_i(z)q_j(z)}{\bar q_G(z)}.
\end{equation}

Let $\vec E$ contain both directions of every undirected edge.
The graph-local inverse-mass affinity gives
\begin{align}
\operatorname{NAssoc}_{\mathrm{soft}}(G)
&=\frac1{\operatorname{vol}(G)}\sum_{(i,j)\in\vec E}W_{ij}a_{\bar q_G}(i,j)\\
&=\sum_{z\in\mathcal Z_G^+}
\frac{\sum_{(i,j)\in\vec E}W_{ij}q_i(z)q_j(z)}{\sum_i d_iq_i(z)}.
\label{eq:soft-nassoc-affinity}
\end{align}
The second expression follows by substituting \cref{eq:graph-aggregate-code}, showing that the graph-local inverse-mass affinity recovers a familiar soft normalized-association surrogate.
To obtain the corresponding fixed-$K$ normalized-cut criterion, we must also specify how empty classes are treated.
We define
\begin{align}
\operatorname{Ncut}_{K,\mathrm{soft}}(G)
&=K-\operatorname{NAssoc}_{\mathrm{soft}}(G)\\
&=K-|\mathcal Z_G^+|
+\sum_{z\in\mathcal Z_G^+}
\frac{\sum_{(i,j)\in\vec E}W_{ij}q_i(z)(1-q_j(z))}
{\sum_i d_iq_i(z)}.
\label{eq:soft-ncut-relational-distortion}
\end{align}
The correction $K-|\mathcal Z_G^+|$ is zero when every soft class is active.
For strictly positive soft assignments every class is active; the correction extends the criterion to hard or limiting assignments with empty classes.

\begin{proposition}[Hard normalized-cut specialization]
\label{prop:normalized-cut}
For a hard assignment $c$ with $K$ positive-volume classes $V_z:=\{i:c(i)=z\}$, \cref{eq:soft-nassoc-affinity} becomes
\begin{equation}
\operatorname{NAssoc}(G)=\sum_{z=1}^K
\frac{\operatorname{assoc}(V_z,V_z)}{\operatorname{vol}(V_z)},
\qquad
\operatorname{Ncut}(G)=K-\operatorname{NAssoc}(G).
\label{eq:hard-ncut}
\end{equation}
These are classical normalized association and normalized cut.
Under the fixed-$K$ extension, an empty class contributes zero association and a unit penalty to normalized cut.
\end{proposition}

\begin{proof}
One-hot membership makes the numerator for class $z$ its directed internal edge weight and the denominator its volume.
The identity $\operatorname{cut}(V_z,V\setminus V_z)=\operatorname{vol}(V_z)-\operatorname{assoc}(V_z,V_z)$ gives the cut form.
The empty-class extension follows from the definition $K-\operatorname{NAssoc}$.
\end{proof}

The hard objective and its clustering correspondences are established~\citep{shi2000normalized,dhillon2004kernel}.
The soft surrogate is also used by neural partitioners such as GAP and the method of Gatti et al.~\citep{nazi2019gap,gatti2022deep}.
Here these established objectives are positioned within the common description of source fidelity and codeword affinity.

The soft ratio uses a deterministic soft mass in its denominator, whereas the volume of a randomly sampled hard class is itself random; it therefore differs from the expectation of hard normalized cut under sampled assignments.
The distinction between ratios of expectations and expected ratios motivates complementary probabilistic-cut formulations~\citep{ghriss2026beyond}.

\subsection{Common structure of the exact correspondences}

The centroid and normalized-association correspondences expose a common representation-side structure.
\begin{equation}
\begin{aligned}
D_{\mathrm{cent}}
&=\frac1{2N^2}\sum_{i,j}\|x_i-x_j\|^2a_{\hat q_Z}(x_i,x_j),\\
\operatorname{NAssoc}_{\mathrm{soft}}
&=\frac1{\operatorname{vol}(G)}\sum_{(i,j)\in\vec E}W_{ij}a_{\bar q_G}(i,j).
\end{aligned}
\label{eq:vq-ncut-relational-correspondence}
\end{equation}
Both use an endogenous inverse-aggregate-mass affinity, while the source-side weighting determines what agreement means: squared Euclidean distance penalizes grouping distant observations, whereas graph edge weight rewards grouping strongly linked vertices.
Thus the same inverse-aggregate-mass affinity construction can realize distinct relational fidelities without identifying their source geometries.
